\section{Spectral phase}\label{sec:phase}

The planar argument will associate unitary operators to the boundary circles
of a surface.  We need a nonnegative numerical invariant that detects the
identity, is subadditive under multiplication, and can be calculated on a
weighted cyclic shift.  We develop these properties in the operator algebra
of an arbitrary group, following the spectral-phase argument of
\cite[Lemma~2.1 and Section~4.3]{OpenAIKervaire2026}.

\subsection{The trace and the choice of phase}

Let $G$ be any group, and let $\{\delta_x:x\in G\}$ be the standard
orthonormal basis of $\ell^2(G)$.  Define left and right translations by
\[
 L(a)\delta_x=\delta_{ax},\qquad
 R_a\delta_x=\delta_{xa^{-1}}.
\]
They commute.  We work in the von Neumann algebra
\[
 \mathcal M=\{T\in\mathcal B(\ell^2(G)):TR_a=R_aT
             \text{ for every }a\in G\},
\]
which contains every $L(a)$.  For a positive integer $d$, regard
$M_d(\mathcal M)$ as operators on $\mathbb C^d\otimes\ell^2(G)$ and set
\begin{equation}\label{eq:group-trace}
 \tau_d(T)=\sum_{j=1}^d(T_{jj})_{1,1},\qquad \tau=\tau_1.
\end{equation}
Here $(T_{jj})_{1,1}$ is the coefficient of $\delta_1$ in
$T_{jj}\delta_1$.  The trace is \emph{unnormalized}: $\tau_d(I)=d$.
For the group von Neumann algebra and this trace construction, see
\cite[Section~3.3]{Jones2009}.

We record why this is a faithful normal trace, including when $G$ is
uncountable.  If $T\in\mathcal M$ and $T\delta_1=\sum_x t_x\delta_x$,
commutation with right translations gives
\[
 T_{x,y}=t_{xy^{-1}}.
\]
For $A,B\in\mathcal M$ with corresponding coefficient families $a,b$,
\[
 \tau(AB)=\sum_{x\in G}a_{x^{-1}}b_x
         =\sum_{x\in G}b_{x^{-1}}a_x=\tau(BA).
\]
These sums converge absolutely by Cauchy--Schwarz; each square-summable
family has countable support.  Summing over the matrix indices proves
traciality of $\tau_d$.  Positivity and normality follow from its expression
as a finite sum of vector functionals.  If $T\geq0$ and $\tau_d(T)=0$, then
\[
 0=\sum_{j=1}^d
   \bigl\|T^{1/2}(e_j\otimes\delta_1)\bigr\|^2.
\]
The square root commutes with all simultaneous right translations, so it
kills $e_j\otimes\delta_x$ for every $j,x$.  Thus $T=0$, proving
faithfulness.

The bounded Borel functional calculus of a normal operator in
$M_d(\mathcal M)$ remains in this algebra.  For a unitary $U$, its spectral
projections define a finite measure
$\mu_U(E)=\tau_d(\mathbf1_E(U))$ on the unit circle $\mathbb T$, and
\begin{equation}\label{eq:phase-spectral-integral}
 \tau_d(f(U))=\int_{\mathbb T}f\,d\mu_U,
 \qquad \mu_U(\mathbb T)=d.
\end{equation}
In particular, uniformly bounded pointwise convergence of Borel functions
implies convergence of these trace evaluations.

Choose the phase with its cut at $1$ by
\begin{equation}\label{eq:phase-definition}
 h(e^{2\pi i t})=t-\lfloor t\rfloor,\qquad
 \ell(U)=\tau_d(h(U)).
\end{equation}
Thus $h(1)=0$, and $h$ is right-continuous in the angle.  Since
$h(U)\geq0$ and $U=\exp(2\pi i h(U))$, faithfulness gives
\begin{equation}\label{eq:phase-faithful}
 \ell(U)\geq0,\qquad \ell(U)=0\ \Longleftrightarrow\ U=I.
\end{equation}
The phase is invariant under unitary conjugation and additive on direct
sums.  On scalar matrices, viewed as matrices over $\mathcal M$, it is the
sum of the phases of the eigenvalues, counted with multiplicity.

\subsection{Subadditivity}

\begin{lemma}[Spectral-phase subadditivity]\label{lem:phase-subadditivity}
For every group $G$, every positive integer $d$, and all unitaries
$U,V\in M_d(\mathcal M)$,
\[
 \ell(UV)\leq\ell(U)+\ell(V).
\]
\end{lemma}

\begin{proof}
Put $D=h(V)$ and join $U$ to $UV$ by the unitary path
\[
 W_u=Ue^{2\pi i uD},\qquad 0\leq u\leq1.
\]
We first calculate the change of a smooth function along this path, and
then approximate the discontinuous function $h$ from the correct side of
its cut.  For $f\in C^\infty(\mathbb T)$ write
$f'(e^{2\pi it})=\frac{d}{dt}f(e^{2\pi it})$.  We claim that
\begin{equation}\label{eq:phase-derivative}
 \frac{d}{du}\tau_d(f(W_u))=\tau_d(f'(W_u)D).
\end{equation}
Indeed, $W'_u=2\pi iW_uD$.  The product and inverse rules, followed by
cyclic permutation inside the trace, give for every integer $k$
\[
 \frac{d}{du}\tau_d(W_u^k)=2\pi i k\,\tau_d(W_u^kD).
\]
This proves \eqref{eq:phase-derivative} for Laurent polynomials.  For smooth
$f$, its Fourier coefficients satisfy
$\sum_{k\in\mathbb Z}(1+|k|)|\widehat f(k)|<\infty$, while
$\|\frac{d}{du}W_u^k\|\leq2\pi|k|\|D\|$.  Thus the Fourier series and
its differentiated series converge uniformly in operator norm, proving
\eqref{eq:phase-derivative} in general.

If $f$ is real-valued and $f'\leq1$, then $I-f'(W_u)\geq0$.  Although
$D$ need not commute with $W_u$, traciality gives the required inequality:
\[
 \tau_d(D)-\tau_d(f'(W_u)D)
 =\tau_d\bigl(D^{1/2}(I-f'(W_u))D^{1/2}\bigr)\geq0.
\]
Integration of \eqref{eq:phase-derivative} therefore yields
\begin{equation}\label{eq:phase-smooth-bound}
 \tau_d(f(UV))-\tau_d(f(U))\leq\tau_d(D).
\end{equation}

For $0<\varepsilon<1$, choose a nonnegative smooth function
$\rho_\varepsilon$ supported in $(0,\varepsilon)$ with integral one, and
define
\[
 f_\varepsilon(e^{2\pi it})
 =\int_0^\varepsilon
   \rho_\varepsilon(s)h(e^{2\pi i(t+s)})\,ds.
\]
This is a smooth function on the circle, with $0\leq f_\varepsilon\leq1$.
For $\delta>0$ the fractional-part function satisfies
\[
 h(e^{2\pi i(t+\delta)})-h(e^{2\pi it})\leq\delta.
\]
Averaging this inequality and taking the difference quotient shows that
$f'_\varepsilon\leq1$.  Right-continuity gives
$f_\varepsilon(\lambda)\to h(\lambda)$ for every $\lambda\in\mathbb T$.
At the cut this follows explicitly from
$0\leq f_\varepsilon(1)\leq\varepsilon$; symmetric smoothing would not
have this property.  Apply \eqref{eq:phase-smooth-bound} to
$f_\varepsilon$ and pass to the limit by
\eqref{eq:phase-spectral-integral}.  Since $\tau_d(D)=\ell(V)$, the
result is the stated inequality.
\end{proof}

The same right-continuity also gives, for any fixed $c\geq0$,
\begin{equation}\label{eq:phase-right-limit}
 \lim_{\varepsilon\downarrow0}
 \ell(e^{2\pi i c\varepsilon}U)=\ell(U).
\end{equation}
In the opposite direction the identity has a jump:
$\ell(e^{-2\pi i\varepsilon}I_d)=d(1-\varepsilon)$ for
$0<\varepsilon<1$.  These two behaviors at the cut will distinguish the
possibly nontrivial boundary circle from the trivial ones in the planar
argument.

\subsection{The phase of a weighted cycle}

Subadditivity will give an upper bound on the operator associated to a
surface.  The following exact calculation will express its phase in terms
of the boundary words.

\begin{lemma}[Weighted-cycle formula]\label{lem:cycle-phase}
Let $a$ be a positive integer, let $Q_1,\ldots,Q_a\in\mathcal M$ be unitary, and define
$T\in M_a(\mathcal M)$ by
\[
 T(e_t\otimes v)=e_{t+1}\otimes Q_tv\quad(1\leq t\leq a),
 \qquad e_{a+1}=e_1.
\]
If $C=Q_aQ_{a-1}\cdots Q_1$ is the product around the cycle starting at
the first coordinate, then
\begin{equation}\label{eq:cycle-phase}
 \ell(T)=\frac{a-1}{2}+\tau(h(C)).
\end{equation}
\end{lemma}

\begin{proof}
Put $\zeta=e^{2\pi i/a}$.  Conjugating $T$ by the scalar diagonal matrix
with entries $1,\zeta,\ldots,\zeta^{a-1}$ gives $\zeta T$.  Hence
$\ell(\zeta^jT)=\ell(T)$ for every integer $j$.  The scalar identity
\[
 \sum_{j=0}^{a-1}h(\zeta^j\lambda)
 =\frac{a-1}{2}+h(\lambda^a),\qquad \lambda\in\mathbb T,
\]
follows by listing the fractional parts of the $a$ equally spaced angles;
it holds also when one of them is zero.  Applying bounded Borel functional
calculus and taking the unnormalized trace gives
\[
 a\ell(T)=\frac{a(a-1)}2+\tau_a(h(T^a)).
\]
The operator $T^a$ is diagonal.  Its first diagonal entry is $C$, and the
other entries are the products around the same cycle with different
starting coordinates.  Consecutive such products are conjugate by the
corresponding $Q_t$, so every diagonal entry has phase trace
$\tau(h(C))$.  Thus $\tau_a(h(T^a))=a\tau(h(C))$, proving
\eqref{eq:cycle-phase}.
\end{proof}
